% TOP_PROBLEM: {"id": 30006899, "title": "Pugh–Shub Stable Ergodicity (C^r Density in Partial Hyperbolicity)", "category_id": 11, "category": {"id": 11, "name": "dynamical_systems", "display_name": "Dynamical Systems", "description": "Problems about long-term behavior of deterministic systems, Hamiltonian dynamics, and periodic orbits.", "slug": "dynamical-systems", "order_index": 11, "created_at": "2026-07-31T15:26:25.671Z"}, "status": "open"}
% FIRST_POSTED: 2026-09-25
% ENTRY_KIND: statement_only
% !TeX program = lualatex
% Definition and sources only; this file is not an LLM solution attempt.
\documentclass[11pt]{article}
\usepackage[margin=25mm]{geometry}
\usepackage{fontspec}
\setmainfont{Latin Modern Roman}
\usepackage{amsmath,amssymb,mathtools}
\usepackage{unicode-math}
\setmathfont{Latin Modern Math}
\usepackage{xurl,hyperref}
\hypersetup{colorlinks=true,urlcolor=blue}
\setcounter{secnumdepth}{0}
\setlength{\parindent}{0pt}
\setlength{\parskip}{5pt}
\setlength{\emergencystretch}{3em}
\begin{document}
\section*{\texorpdfstring{Pugh–Shub Stable Ergodicity (C\textasciicircum{}r Density in Partial Hyperbolicity)}{Pugh–Shub Stable Ergodicity (C\textasciicircum{}r Density in Partial Hyperbolicity)}}\label{30006899}
\subsection{Definitions and mathematical statement}
Let $(M,m)$ be a compact connected boundaryless smooth manifold with smooth probability volume. Partial hyperbolicity is a continuous invariant splitting $TM=E^u\oplus E^c\oplus E^s$ with nonzero $E^u,E^s$, possible zero center, and a common iterate $N$ such that unstable unit vectors expand by more than two, stable ones contract below one-half, and earlier nonzero bundles dominate later ones by a factor greater than two at the same base point. A map $g$ is stably ergodic here if a $C^1$ neighborhood of it has every $C^2$ volume-preserving member ergodic for $m$, meaning every invariant measurable set has measure zero or one. For every $r\in[2,\infty]$, the assertion is density of these $g$ in $\operatorname{PH}_m^r(M)$ in the relative $C^r$ topology. Nonintegral finite $r$ uses H\"older regularity. The density topology and the stability topology are intentionally different.
\par\smallskip\textit{Formulation and concept references:} \href{https://arxiv.org/pdf/1709.04983}{[S1]}.

\subsection{Short English statement}
Can every volume-preserving partially hyperbolic map be approximated in its full $C^r$ topology by one whose nearby $C^2$ conservative maps remain ergodic?
\subsection{Sources}
\begin{itemize}
\item[S1]Pugh–Shub Stable Ergodicity (C\textasciicircum{}r Density in Partial Hyperbolicity); exact editorial selection using the primary formulations and hypotheses recorded in the source receipts.
\url{https://arxiv.org/pdf/1709.04983}.
\textit{Formulation source.}
\end{itemize}
\end{document}
